\documentclass[aspectratio=169]{beamer}

\usetheme{Madrid}
\usepackage{booktabs}
\usepackage{xcolor}
\setbeamertemplate{navigation symbols}{}

\title[Applying and hardening the LULC classifier]{Applying the LULC classifier on real sites: robustness, data adequacy, and validation}
\author[Salil Gujar]{Salil Gujar \\ \footnotesize M.Tech CSE \quad Entry No.\ 2025MCS2106}
\date{July 2026}

\begin{document}

\frame{\titlepage}

% ------------------------------------------------------------------
\begin{frame}{Where we are, and what was advised}
\begin{columns}[T]
\begin{column}{0.5\textwidth}
\emph{Done so far}
\begin{itemize}\small
  \item A 4-class base map over India at 10\,m, with a living hierarchy: split or add a class,
  give a few example polygons, and retrain that node on the fly.
  \item A git-backed zoo of Model and Dataset cards, with save and reload, merging, base and
  year selection, and publishing to a shared repository.
\end{itemize}
\end{column}
\begin{column}{0.5\textwidth}
\emph{What was advised}
\begin{itemize}\small
  \item Start applying the framework on real sites: acacia against non-acacia, tea against
  non-tea, and mining.
  \item Make the models robust across years by training on more than one year.
  \item Judge training data by how much of the area it covers and how it is spread, not by a
  raw count.
  \item Validate an uploaded scheme before running it, and revisit the data choices: inference
  year and Tessera.
\end{itemize}
\end{column}
\end{columns}
\end{frame}

% ------------------------------------------------------------------
\begin{frame}{Standard testing sites}
\begin{itemize}
  \item We fixed a home strip around IIT Delhi and Sanjay Van: built-up, tree cover, and a small
  water body in one small area, with acacia and non-acacia crowns inside it.
  \item Three such sites are now default areas in the tool, chosen from a menu without typing
  coordinates: IIT Delhi with Sanjay Van, Asola Bhatti, and Jalpaiguri, alongside the Assam tea
  belt already present.
  \item For acacia, we now hold 912 confidently labelled tree crowns, 336 acacia and 576
  non-acacia, prepared as ready training examples for the tool.
  \item Asola Bhatti serves as a false-positive probe: old mines there were reclaimed into
  built-up and scrub, so a mining model should stay quiet and an acacia model may not.
\end{itemize}
\end{frame}

% ------------------------------------------------------------------
\begin{frame}{Robustness across years}
\begin{itemize}
  \item Vegetation changes from year to year, and Alpha Earth has a separate yearly image, so a
  model trained on one year can drift on another.
  \item A split can now be trained on several years at once. The same polygon sampled in
  different years becomes more training examples, while a whole polygon is still held out
  together, so the test stays honest.
  \item On acacia and non-acacia, training on 2019, 2021, and 2023 and testing on the unseen
  2020 and 2024 raised accuracy from 0.635 to 0.745, a gain of about eleven points over a
  single-year model.
  \item On the four base classes the multi-year and single-year models are the same, both near
  0.89 on the unseen years. The coarse classes look alike year to year, so a single year already
  generalises; the gain from multiple years shows up on the finer, year-sensitive splits.
\end{itemize}
\end{frame}

% ------------------------------------------------------------------
\begin{frame}{Testing on the sites: what we found}
\begin{itemize}\small
  \item Tea against non-tea separates cleanly. Acacia against non-acacia is genuinely hard,
  since it is a species distinction, but it is where training on several years helps most.
  \item The mining model detects real mines well: an active coalfield (Jharia) reads 71\% mining.
  On reclaimed, mine-like ground it raises false positives: reclaimed Asola reads 17\%, which is
  exactly the concern raised for Asola Bhatti. Same detector on both, so the two are comparable.
\end{itemize}
\vspace{0.5em}
\centering
\begin{tabular}{ll}
\toprule
Test & Result \\
\midrule
Tea vs non-tea & 0.957 held-out accuracy \\
Acacia vs non-acacia & 0.745 on unseen years (multi-year); a hard split \\
Mining detection & 0.859 accuracy, 0.854 recall on real mines \\
Mining false positives & 0.139 on non-mining ground truth \\
Jharia, active coalfield & 71.2\% of the area flagged mining (real detections) \\
Asola Bhatti, reclaimed & 17.1\% flagged mining (mostly false positives) \\
\bottomrule
\end{tabular}
\end{frame}

% ------------------------------------------------------------------
\begin{frame}{Judging whether there is enough training data}
\begin{itemize}
  \item A fixed number of labelled points means little on its own: the same count is plenty over
  a village and almost nothing over a district.
  \item We now report coverage: how much of the area about to be classified actually falls inside
  a labelled polygon. It is read against the size of that area, so it rises on a small area and
  falls on a large one.
  \item Alongside it we keep the spread measure, the spatial evenness of the labels across the
  region, since data clustered in one place skews a model.
\end{itemize}
\end{frame}

% ------------------------------------------------------------------
\begin{frame}{Validating a scheme before it runs}
\begin{itemize}
  \item A user can upload a saved scheme: the class hierarchy and the ordered steps that built
  it. Until now it was applied first and any problem was noticed afterwards.
  \item There is now a check that runs before anything is applied and changes nothing. It reports
  errors that block loading and warnings that do not.
  \item It confirms the hierarchy is well formed, that every recorded operation is one the tool
  knows and carries its required fields, and that each classifier the scheme refers to is present.
  \item Loading now runs this check first, so a broken file is refused before it can alter the
  live tree.
\end{itemize}
\end{frame}

% ------------------------------------------------------------------
\begin{frame}{Choosing the data: years and Tessera}
\begin{itemize}
  \item The inference-year picker now carries a note on how far a ground truth stays reliable
  across years, so the user knows what a chosen year buys.
  \item Tessera remains available as the second feature source in Detailed mode. It has usable
  India coverage only for 2024, and the note states that other years can be requested.
\end{itemize}
\end{frame}

% ------------------------------------------------------------------
\begin{frame}
\centering
\Large Thank you
\end{frame}

\end{document}
